\ifdefined\IntervalProjectMaster
\def\IntervalProjectEnd{}
\section{Explicit multi-interval constructions}
\label{sec:explicit-m4-m5}
\label{sec:explicit-constructions}
\else
\documentclass{article}
\usepackage[utf8]{inputenc}
\usepackage{amsmath}
\usepackage{amssymb}
\usepackage{amsthm}
\usepackage{parskip}
\usepackage{url}

\newtheorem{proposition}{Proposition}
\newcommand{\code}[1]{\path{#1}}

\title{Multi-interval examples improving the one-interval bound}
\author{Ferran Espu\~na}
\date{\today}
\def\IntervalProjectEnd{\end{document}}

\begin{document}

\maketitle
\fi

We write \(\mathbb T=\mathbb R/\mathbb Z\). A set \(A\subset \mathbb T\) is
\(m\)-sum-free if
\[
    0\notin A+A-mA.
\]
Equivalently, there are no \(x,y,z\in A\) with \(x+y=mz\) in \(\mathbb T\).

For comparison, suppose \(A=I\) is a single open interval of length \(\alpha\).
Then \(I+I-mI\) is an open interval of length \((m+2)\alpha\) on the real line.
If \((m+2)\alpha>1\), its image in \(\mathbb T\) is all of \(\mathbb T\). Thus
the trivial one-interval bound is
\[
    \alpha\leq \frac1{m+2}.
\]
The examples below exceed this bound for \(m=4,5,8,9,12,13\). We also record
their exact endpoint certificates.  They are concrete examples of the two
preceding phenomena: the \(m=5\) example realizes the equally spaced lower
bound, the \(m=4,8,9\) examples add one short interval to that baseline, and
the \(m=12\) and \(m=13\) examples add three and two short intervals,
respectively.

We will use the following finite certificate. If
\[
    I_r=(a_r,b_r),
\]
then
\[
    I_i+I_j-mI_\ell
    =
    \bigl(a_i+a_j-mb_\ell,\; b_i+b_j-ma_\ell\bigr).
\]
It is therefore enough, for each \(i\leq j\) and each \(\ell\), to find an
integer \(n\) such that
\[
    n\leq a_i+a_j-mb_\ell
    \quad\text{and}\quad
    b_i+b_j-ma_\ell\leq n+1.
\]
Then the open interval \(I_i+I_j-mI_\ell\) contains no integer, and hence
contributes no point equal to \(0\) in \(\mathbb T\).

\begin{proposition}[A \(4\)-sum-free set of three open intervals]
Let
\[
    A_4=
    \left(\frac13,\frac5{12}\right)
    \cup
    \left(\frac{25}{48},\frac{101}{192}\right)
    \cup
    \left(\frac7{12},\frac23\right)
    \subset \mathbb T.
\]
Then \(A_4\) is \(4\)-sum-free and
\[
    \mu(A_4)
    =
    \frac1{12}+\frac1{192}+\frac1{12}
    =
    \frac{11}{64}
    >
    \frac16.
\]
Thus the one-interval bound for \(m=4\) is improved by \(1/192\).
\end{proposition}

\begin{proof}
Let the three intervals in \(A_4\) be \(I_1,I_2,I_3\), in the displayed order.
For each triple \(i\leq j\) and \(\ell\), the following table gives
\[
    L_{ij\ell}=a_i+a_j-4b_\ell,
    \qquad
    R_{ij\ell}=b_i+b_j-4a_\ell,
\]
and an integer \(n\) with \(n\leq L_{ij\ell}<R_{ij\ell}\leq n+1\).
\[
\begin{array}{c|c|c|c}
(i,j,\ell) & n & L_{ij\ell} & R_{ij\ell} \\
\hline
(1,1,1) & -1 & -1 & -\frac12 \\
(1,1,2) & -2 & -\frac{23}{16} & -\frac54 \\
(1,1,3) & -2 & -2 & -\frac32 \\
(1,2,1) & -1 & -\frac{13}{16} & -\frac{25}{64} \\
(1,2,2) & -2 & -\frac54 & -\frac{73}{64} \\
(1,2,3) & -2 & -\frac{29}{16} & -\frac{89}{64} \\
(1,3,1) & -1 & -\frac34 & -\frac14 \\
(1,3,2) & -2 & -\frac{19}{16} & -1 \\
(1,3,3) & -2 & -\frac74 & -\frac54 \\
(2,2,1) & -1 & -\frac58 & -\frac9{32} \\
(2,2,2) & -2 & -\frac{17}{16} & -\frac{33}{32} \\
(2,2,3) & -2 & -\frac{13}{8} & -\frac{41}{32} \\
(2,3,1) & -1 & -\frac9{16} & -\frac9{64} \\
(2,3,2) & -1 & -1 & -\frac{57}{64} \\
(2,3,3) & -2 & -\frac{25}{16} & -\frac{73}{64} \\
(3,3,1) & -1 & -\frac12 & 0 \\
(3,3,2) & -1 & -\frac{15}{16} & -\frac34 \\
(3,3,3) & -2 & -\frac32 & -1
\end{array}
\]
Since the intervals are open, equality with an integer endpoint in the table
does not introduce an integer point. Hence every interval
\(I_i+I_j-4I_\ell\) misses all integers, so \(0\notin A_4+A_4-4A_4\) in
\(\mathbb T\).
\end{proof}

\begin{proposition}[A \(5\)-sum-free set of two open intervals]
Let
\[
    A_5=
    \left(\frac2{35},\frac17\right)
    \cup
    \left(\frac67,\frac{33}{35}\right)
    \subset \mathbb T.
\]
Then \(A_5\) is \(5\)-sum-free and
\[
    \mu(A_5)
    =
    \frac3{35}+\frac3{35}
    =
    \frac6{35}
    >
    \frac17.
\]
Thus the one-interval bound for \(m=5\) is improved by \(1/35\).
\end{proposition}

\begin{proof}
The numerical experiment found a third interval of length \(0\), namely
\((6/7,6/7)\). We omit it here and keep only the two non-empty intervals above.
Let these intervals be \(I_1,I_2\).

For each triple \(i\leq j\) and \(\ell\), set
\[
    L_{ij\ell}=a_i+a_j-5b_\ell,
    \qquad
    R_{ij\ell}=b_i+b_j-5a_\ell.
\]
The following table gives an integer \(n\) with
\[
    n\leq L_{ij\ell}<R_{ij\ell}\leq n+1.
\]
\[
\begin{array}{c|c|c|c}
(i,j,\ell) & n & L_{ij\ell} & R_{ij\ell} \\
\hline
(1,1,1) & -1 & -\frac35 & 0 \\
(1,1,2) & -5 & -\frac{23}{5} & -4 \\
(1,2,1) & 0 & \frac15 & \frac45 \\
(1,2,2) & -4 & -\frac{19}{5} & -\frac{16}{5} \\
(2,2,1) & 1 & 1 & \frac85 \\
(2,2,2) & -3 & -3 & -\frac{12}{5}
\end{array}
\]
Again, the intervals are open, so the integer endpoints shown in the table are
not included. Therefore every interval \(I_i+I_j-5I_\ell\) misses all integers,
and \(0\notin A_5+A_5-5A_5\) in \(\mathbb T\).
\end{proof}

\begin{proposition}[Exact constructions recovered for \(m=8,9\)]
The open sets
\[
\begin{split}
    A_8={}&
    \left(\frac7{120},\frac{13}{120}\right)
    \cup\left(\frac{29}{240},\frac{119}{960}\right)
    \cup\left(\frac{11}{60},\frac7{30}\right)
    \cup\left(\frac{14}{15},\frac{59}{60}\right),\\
    A_9={}&
    \left(\frac{39}{77},\frac{386}{693}\right)
    \cup\left(\frac{428}{693},\frac{463}{693}\right)
    \cup\left(\frac{4237}{6237},\frac{4244}{6237}\right)
    \cup\left(\frac{505}{693},\frac{60}{77}\right)
\end{split}
\]
are respectively \(8\)-sum-free and \(9\)-sum-free.  Consequently,
\[
    d_8(\mathbb T)\geq\mu(A_8)=\frac{49}{320},
    \qquad
    d_9(\mathbb T)\geq\mu(A_9)=\frac{136}{891}.
\]
\end{proposition}

\begin{proof}
Number the four intervals in each set in their displayed order.  For each
pair \((i,j)\), the four entries in the corresponding row below give the
integer \(n\) for \(\ell=1,2,3,4\) in the finite certificate above.
\[
\begin{array}{c|rrrr}
\multicolumn{5}{c}{m=8}\\
(i,j)\backslash\ell&1&2&3&4\\ \hline
(1,1)&-1&-1&-2&-8\\
(1,2)&-1&-1&-2&-8\\
(1,3)&-1&-1&-2&-8\\
(1,4)& 0& 0&-1&-7\\
(2,2)&-1&-1&-2&-8\\
(2,3)&-1&-1&-2&-8\\
(2,4)& 0& 0&-1&-7\\
(3,3)&-1&-1&-2&-8\\
(3,4)& 0& 0&-1&-7\\
(4,4)& 1& 0& 0&-6
\end{array}
\qquad
\begin{array}{c|rrrr}
\multicolumn{5}{c}{m=9}\\
(i,j)\backslash\ell&1&2&3&4\\ \hline
(1,1)&-4&-5&-6&-6\\
(1,2)&-4&-5&-5&-6\\
(1,3)&-4&-5&-5&-6\\
(1,4)&-4&-5&-5&-6\\
(2,2)&-4&-5&-5&-6\\
(2,3)&-4&-5&-5&-6\\
(2,4)&-4&-5&-5&-6\\
(3,3)&-4&-5&-5&-6\\
(3,4)&-4&-5&-5&-6\\
(4,4)&-4&-5&-5&-6
\end{array}
\]
Exact substitution of the displayed endpoints gives
\[
    n\leq a_i+a_j-mb_\ell
    <b_i+b_j-ma_\ell\leq n+1
\]
in all 40 cases for each multiplier.  The smallest nonzero slack is
\(1/16\) for \(A_8\) and \(5/81\) for \(A_9\); equality in the remaining
tight sides is harmless because all intervals are open.  Thus no lifted
interaction contains an integer.  Finally, the three long components have
total length \(3/20\) for \(A_8\) and \(5/33\) for \(A_9\), while the short
components have lengths \(1/320\) and \(1/891\), respectively, giving the
stated measures.
\end{proof}

\begin{proposition}[Exact constructions with several added intervals]
The open sets
\[
\begin{split}
 A_{12}={}&
 \left(\frac{17}{420},\frac8{105}\right)
 \cup\left(\frac{23}{280},\frac{419}{5040}\right)
 \cup\left(\frac{7103}{60480},\frac{33}{280}\right)
 \cup\left(\frac{13}{105},\frac{67}{420}\right)\\
 &\cup\left(\frac{139}{840},\frac{839}{5040}\right)
 \cup\left(\frac{29}{140},\frac{17}{70}\right)
 \cup\left(\frac{67}{70},\frac{139}{140}\right),\\
 A_{13}={}&
 \left(\frac{632}{27885},\frac{643}{27885}\right)
 \cup\left(\frac{797}{27885},\frac{808}{27885}\right)
 \cup\left(\frac{74}{2145},\frac{151}{2145}\right)\\
 &\cup\left(\frac{239}{2145},\frac{316}{2145}\right)
 \cup\left(\frac{404}{2145},\frac{481}{2145}\right)
 \cup\left(\frac{2054}{2145},\frac{2131}{2145}\right)
\end{split}
\]
are respectively \(12\)-sum-free and \(13\)-sum-free.  In particular,
\[
    d_{12}(\mathbb T)\geq\frac{251}{1728},
    \qquad
    d_{13}(\mathbb T)\geq\frac{122}{845}.
\]
\end{proposition}

\begin{proof}
For each set, order the displayed intervals as written and, for every
\(i\leq j\) and \(\ell\), take
\[
    n=\left\lfloor a_i+a_j-mb_\ell\right\rfloor.
\]
For \(A_{12}\), direct fraction arithmetic in all 196 cases gives
\[
    n\leq a_i+a_j-mb_\ell
    <b_i+b_j-ma_\ell\leq n+1.
\]
There are 14 tight endpoint inequalities, and the minimum positive endpoint
slack is \(1/144\).  For \(A_{13}\), the same calculation in all 126 cases
has 13 tight endpoint inequalities and minimum positive slack \(7/169\).
The finite certificate therefore proves sum-freeness in both cases.

The four base components of \(A_{12}\) contribute \(1/7\); two added
components have length \(1/1008\), and the third has length \(5/12096\).
Consequently,
\[
    \mu(A_{12})
    =\frac17+\frac2{1008}+\frac5{12096}
    =\frac{251}{1728}.
\]
Similarly, the four base components of \(A_{13}\) contribute \(28/195\),
and the other two have length \(1/2535\) each, giving
\(\mu(A_{13})=122/845\).  Both exact calculations are reproduced in the
directory \code{tiny_intervals_three_extra_exact/} by the script
\code{verify_exact_three_extra.py}.
\end{proof}

\IntervalProjectEnd
